\section{Sort a Stack using Recursion}
\label{sec:rec-sort-stack}

\problemstatement{Given a stack, sort it in increasing order (so the
smallest element ends up on top) using \textbf{only} \textit{push} and
\textit{pop}, and recursion -- no second stack, array, or queue allowed.}

\begin{patternbox}
This is Section~\ref{sec:rec-reverse-stack}'s ``insert at the bottom''
helper, generalized to ``insert at the \emph{correct sorted position}.''
The overall structure -- pop the top, recursively solve the rest, then
reinsert the popped value using a specialized helper -- is identical;
only the helper's insertion rule changes from ``always go to the
bottom'' to ``go just above everything smaller, and below everything
larger.''
\end{patternbox}

\subsection*{Understanding the problem carefully}

If the remainder of the stack (everything below the current top) is
already correctly sorted, inserting the popped-off top element into its
correct position relative to that sorted remainder finishes the job --
this is exactly insertion sort's core idea, expressed with a stack
playing the role that an array's shifting-elements-right step normally
plays.

\begin{ideabox}[Insert into a Sorted Stack at the Correct Position]
$\textsc{InsertSorted}(\textit{stack}, x)$: if $\textit{stack}$ is empty
\emph{or} its top is $\le x$, push $x$ and return (its correct position
has been found: everything remaining below is $\le x$, satisfying sorted
order with the top of the stack as the largest visible value at this
point... more precisely, the stack is kept increasing from top to
bottom, so $x$ belongs here once it is at least as large as what is
currently on top). Otherwise, pop the top into $\textit{temp}$,
recursively call $\textsc{InsertSorted}(\textit{stack}, x)$ (continue
searching deeper for $x$'s correct position), then push $\textit{temp}$
back on top.
\end{ideabox}

\subsection*{Why this insertion rule keeps the stack sorted at every step}

Assume, inductively, that before inserting $x$, the stack is already
sorted (say, increasing from top to bottom: the smallest element on top).
The helper pops elements off only while they are \emph{larger} than $x$
-- precisely the elements that must end up \emph{above} $x$ once $x$ is
correctly placed. The moment it finds an element $\le x$ (or the stack
empties), that is exactly the right depth to insert $x$, and pushing the
previously-popped larger elements back afterward restores them directly
on top of $x$, in their original relative order (preserved by the
recursive call stack's own LIFO unwinding) -- so the stack remains
sorted after the insertion completes.

\subsection*{Algorithm}

\begin{enumerate}
  \item $\textsc{InsertSorted}(\textit{stack}, x)$:
  \begin{enumerate}
    \item If $\textit{stack}$ is empty or top $\le x$: push $x$; return.
    \item Pop top into $\textit{temp}$.
    \item $\textsc{InsertSorted}(\textit{stack}, x)$.
    \item Push $\textit{temp}$.
  \end{enumerate}
  \item $\textsc{SortStack}(\textit{stack})$:
  \begin{enumerate}
    \item If $\textit{stack}$ is empty: return.
    \item Pop top into $\textit{temp}$.
    \item $\textsc{SortStack}(\textit{stack})$.
    \item $\textsc{InsertSorted}(\textit{stack}, \textit{temp})$.
  \end{enumerate}
\end{enumerate}

\subsection*{Dry run}

\dryrun{Sort the stack $[3,1,2]$ (top to bottom: $2,1,3$).}

\begin{itemize}
  \item Pop $2$. Recurse on $[3,1]$ (top to bottom: $1,3$).
  \item Pop $1$. Recurse on $[3]$.
  \item Pop $3$. Recurse on $[]$ (base case).
  \item $\textsc{InsertSorted}([], 3)$: empty, push. Stack$=[3]$.
  \item Back in the ``pop $1$'' frame: $\textsc{InsertSorted}([3], 1)$:
        top is $3 > 1$, pop $3$ into \textit{temp}; recurse
        $\textsc{InsertSorted}([], 1)$: empty, push $1$, stack$=[1]$;
        push back \textit{temp} ($3$): stack$=[1,3]$ (top to bottom:
        $1,3$ -- sorted).
  \item Back in the ``pop $2$'' frame: $\textsc{InsertSorted}([1,3],
        2)$: top is $1 \le 2$, push $2$ directly. Stack$=[2,1,3]$ (top
        to bottom: $2,1,3$).
\end{itemize}

Final stack, top to bottom: $1,2,3$ -- correctly sorted, smallest on
top.

\subsection*{C++ implementation}

\begin{lstlisting}
#include <bits/stdc++.h>
using namespace std;

void insertSorted(stack<int>& st, int x) {
    if (st.empty() || st.top() <= x) {
        st.push(x);
        return;
    }
    int temp = st.top();
    st.pop();
    insertSorted(st, x);
    st.push(temp);
}

void sortStack(stack<int>& st) {
    if (st.empty()) return;

    int temp = st.top();
    st.pop();
    sortStack(st);
    insertSorted(st, temp);
}

int main() {
    stack<int> st;
    st.push(3); st.push(1); st.push(2);

    sortStack(st);

    while (!st.empty()) {
        cout << st.top() << " ";
        st.pop();
    }
    cout << endl;
    return 0;
}
\end{lstlisting}

\begin{complexitybox}
\textbf{Time Complexity.} $O(n^2)$ in the worst case -- the same
insertion-sort-style bound as Section~\ref{sec:rec-reverse-stack}, since
$\textsc{InsertSorted}$ can pop and re-push up to $O(k)$ elements for a
stack of size $k$, summed over $n$ insertions.
\textbf{Space Complexity.} $O(n)$ for the recursion stack at the deepest
point.
\end{complexitybox}

\begin{takeawaybox}
Once you have solved one ``insert at a specific position using only
recursion'' problem (insert at the bottom, in
Section~\ref{sec:rec-reverse-stack}), solving a closely related one is
just a matter of changing the stopping condition inside the same
recursive shape: stop unconditionally at the bottom for reversal, stop as
soon as a smaller-or-equal element is found for sorting.
\end{takeawaybox}
